% Options for packages loaded elsewhere
% Options for packages loaded elsewhere
\PassOptionsToPackage{unicode}{hyperref}
\PassOptionsToPackage{hyphens}{url}
\PassOptionsToPackage{dvipsnames,svgnames,x11names}{xcolor}
%
\documentclass[
  letterpaper,
  DIV=11,
  numbers=noendperiod]{scrreprt}
\usepackage{xcolor}
\usepackage{amsmath,amssymb}
\setcounter{secnumdepth}{5}
\usepackage{iftex}
\ifPDFTeX
  \usepackage[T1]{fontenc}
  \usepackage[utf8]{inputenc}
  \usepackage{textcomp} % provide euro and other symbols
\else % if luatex or xetex
  \usepackage{unicode-math} % this also loads fontspec
  \defaultfontfeatures{Scale=MatchLowercase}
  \defaultfontfeatures[\rmfamily]{Ligatures=TeX,Scale=1}
\fi
\usepackage{lmodern}
\ifPDFTeX\else
  % xetex/luatex font selection
\fi
% Use upquote if available, for straight quotes in verbatim environments
\IfFileExists{upquote.sty}{\usepackage{upquote}}{}
\IfFileExists{microtype.sty}{% use microtype if available
  \usepackage[]{microtype}
  \UseMicrotypeSet[protrusion]{basicmath} % disable protrusion for tt fonts
}{}
\makeatletter
\@ifundefined{KOMAClassName}{% if non-KOMA class
  \IfFileExists{parskip.sty}{%
    \usepackage{parskip}
  }{% else
    \setlength{\parindent}{0pt}
    \setlength{\parskip}{6pt plus 2pt minus 1pt}}
}{% if KOMA class
  \KOMAoptions{parskip=half}}
\makeatother
% Make \paragraph and \subparagraph free-standing
\makeatletter
\ifx\paragraph\undefined\else
  \let\oldparagraph\paragraph
  \renewcommand{\paragraph}{
    \@ifstar
      \xxxParagraphStar
      \xxxParagraphNoStar
  }
  \newcommand{\xxxParagraphStar}[1]{\oldparagraph*{#1}\mbox{}}
  \newcommand{\xxxParagraphNoStar}[1]{\oldparagraph{#1}\mbox{}}
\fi
\ifx\subparagraph\undefined\else
  \let\oldsubparagraph\subparagraph
  \renewcommand{\subparagraph}{
    \@ifstar
      \xxxSubParagraphStar
      \xxxSubParagraphNoStar
  }
  \newcommand{\xxxSubParagraphStar}[1]{\oldsubparagraph*{#1}\mbox{}}
  \newcommand{\xxxSubParagraphNoStar}[1]{\oldsubparagraph{#1}\mbox{}}
\fi
\makeatother

\usepackage{color}
\usepackage{fancyvrb}
\newcommand{\VerbBar}{|}
\newcommand{\VERB}{\Verb[commandchars=\\\{\}]}
\DefineVerbatimEnvironment{Highlighting}{Verbatim}{commandchars=\\\{\}}
% Add ',fontsize=\small' for more characters per line
\usepackage{framed}
\definecolor{shadecolor}{RGB}{241,243,245}
\newenvironment{Shaded}{\begin{snugshade}}{\end{snugshade}}
\newcommand{\AlertTok}[1]{\textcolor[rgb]{0.68,0.00,0.00}{#1}}
\newcommand{\AnnotationTok}[1]{\textcolor[rgb]{0.37,0.37,0.37}{#1}}
\newcommand{\AttributeTok}[1]{\textcolor[rgb]{0.40,0.45,0.13}{#1}}
\newcommand{\BaseNTok}[1]{\textcolor[rgb]{0.68,0.00,0.00}{#1}}
\newcommand{\BuiltInTok}[1]{\textcolor[rgb]{0.00,0.23,0.31}{#1}}
\newcommand{\CharTok}[1]{\textcolor[rgb]{0.13,0.47,0.30}{#1}}
\newcommand{\CommentTok}[1]{\textcolor[rgb]{0.37,0.37,0.37}{#1}}
\newcommand{\CommentVarTok}[1]{\textcolor[rgb]{0.37,0.37,0.37}{\textit{#1}}}
\newcommand{\ConstantTok}[1]{\textcolor[rgb]{0.56,0.35,0.01}{#1}}
\newcommand{\ControlFlowTok}[1]{\textcolor[rgb]{0.00,0.23,0.31}{\textbf{#1}}}
\newcommand{\DataTypeTok}[1]{\textcolor[rgb]{0.68,0.00,0.00}{#1}}
\newcommand{\DecValTok}[1]{\textcolor[rgb]{0.68,0.00,0.00}{#1}}
\newcommand{\DocumentationTok}[1]{\textcolor[rgb]{0.37,0.37,0.37}{\textit{#1}}}
\newcommand{\ErrorTok}[1]{\textcolor[rgb]{0.68,0.00,0.00}{#1}}
\newcommand{\ExtensionTok}[1]{\textcolor[rgb]{0.00,0.23,0.31}{#1}}
\newcommand{\FloatTok}[1]{\textcolor[rgb]{0.68,0.00,0.00}{#1}}
\newcommand{\FunctionTok}[1]{\textcolor[rgb]{0.28,0.35,0.67}{#1}}
\newcommand{\ImportTok}[1]{\textcolor[rgb]{0.00,0.46,0.62}{#1}}
\newcommand{\InformationTok}[1]{\textcolor[rgb]{0.37,0.37,0.37}{#1}}
\newcommand{\KeywordTok}[1]{\textcolor[rgb]{0.00,0.23,0.31}{\textbf{#1}}}
\newcommand{\NormalTok}[1]{\textcolor[rgb]{0.00,0.23,0.31}{#1}}
\newcommand{\OperatorTok}[1]{\textcolor[rgb]{0.37,0.37,0.37}{#1}}
\newcommand{\OtherTok}[1]{\textcolor[rgb]{0.00,0.23,0.31}{#1}}
\newcommand{\PreprocessorTok}[1]{\textcolor[rgb]{0.68,0.00,0.00}{#1}}
\newcommand{\RegionMarkerTok}[1]{\textcolor[rgb]{0.00,0.23,0.31}{#1}}
\newcommand{\SpecialCharTok}[1]{\textcolor[rgb]{0.37,0.37,0.37}{#1}}
\newcommand{\SpecialStringTok}[1]{\textcolor[rgb]{0.13,0.47,0.30}{#1}}
\newcommand{\StringTok}[1]{\textcolor[rgb]{0.13,0.47,0.30}{#1}}
\newcommand{\VariableTok}[1]{\textcolor[rgb]{0.07,0.07,0.07}{#1}}
\newcommand{\VerbatimStringTok}[1]{\textcolor[rgb]{0.13,0.47,0.30}{#1}}
\newcommand{\WarningTok}[1]{\textcolor[rgb]{0.37,0.37,0.37}{\textit{#1}}}

\usepackage{longtable,booktabs,array}
\newcounter{none} % for unnumbered tables
\usepackage{calc} % for calculating minipage widths
% Correct order of tables after \paragraph or \subparagraph
\usepackage{etoolbox}
\makeatletter
\patchcmd\longtable{\par}{\if@noskipsec\mbox{}\fi\par}{}{}
\makeatother
% Allow footnotes in longtable head/foot
\IfFileExists{footnotehyper.sty}{\usepackage{footnotehyper}}{\usepackage{footnote}}
\makesavenoteenv{longtable}
\usepackage{graphicx}
\makeatletter
\newsavebox\pandoc@box
\newcommand*\pandocbounded[1]{% scales image to fit in text height/width
  \sbox\pandoc@box{#1}%
  \Gscale@div\@tempa{\textheight}{\dimexpr\ht\pandoc@box+\dp\pandoc@box\relax}%
  \Gscale@div\@tempb{\linewidth}{\wd\pandoc@box}%
  \ifdim\@tempb\p@<\@tempa\p@\let\@tempa\@tempb\fi% select the smaller of both
  \ifdim\@tempa\p@<\p@\scalebox{\@tempa}{\usebox\pandoc@box}%
  \else\usebox{\pandoc@box}%
  \fi%
}
% Set default figure placement to htbp
\def\fps@figure{htbp}
\makeatother


% definitions for citeproc citations
\NewDocumentCommand\citeproctext{}{}
\NewDocumentCommand\citeproc{mm}{%
  \begingroup\def\citeproctext{#2}\cite{#1}\endgroup}
\makeatletter
 % allow citations to break across lines
 \let\@cite@ofmt\@firstofone
 % avoid brackets around text for \cite:
 \def\@biblabel#1{}
 \def\@cite#1#2{{#1\if@tempswa , #2\fi}}
\makeatother
\newlength{\cslhangindent}
\setlength{\cslhangindent}{1.5em}
\newlength{\csllabelwidth}
\setlength{\csllabelwidth}{3em}
\newenvironment{CSLReferences}[2] % #1 hanging-indent, #2 entry-spacing
 {\begin{list}{}{%
  \setlength{\itemindent}{0pt}
  \setlength{\leftmargin}{0pt}
  \setlength{\parsep}{0pt}
  % turn on hanging indent if param 1 is 1
  \ifodd #1
   \setlength{\leftmargin}{\cslhangindent}
   \setlength{\itemindent}{-1\cslhangindent}
  \fi
  % set entry spacing
  \setlength{\itemsep}{#2\baselineskip}}}
 {\end{list}}
\usepackage{calc}
\newcommand{\CSLBlock}[1]{\hfill\break\parbox[t]{\linewidth}{\strut\ignorespaces#1\strut}}
\newcommand{\CSLLeftMargin}[1]{\parbox[t]{\csllabelwidth}{\strut#1\strut}}
\newcommand{\CSLRightInline}[1]{\parbox[t]{\linewidth - \csllabelwidth}{\strut#1\strut}}
\newcommand{\CSLIndent}[1]{\hspace{\cslhangindent}#1}



\setlength{\emergencystretch}{3em} % prevent overfull lines

\providecommand{\tightlist}{%
  \setlength{\itemsep}{0pt}\setlength{\parskip}{0pt}}



 


\KOMAoption{captions}{tableheading}
\makeatletter
\@ifpackageloaded{tcolorbox}{}{\usepackage[skins,breakable]{tcolorbox}}
\@ifpackageloaded{fontawesome5}{}{\usepackage{fontawesome5}}
\definecolor{quarto-callout-color}{HTML}{909090}
\definecolor{quarto-callout-note-color}{HTML}{0758E5}
\definecolor{quarto-callout-important-color}{HTML}{CC1914}
\definecolor{quarto-callout-warning-color}{HTML}{EB9113}
\definecolor{quarto-callout-tip-color}{HTML}{00A047}
\definecolor{quarto-callout-caution-color}{HTML}{FC5300}
\definecolor{quarto-callout-color-frame}{HTML}{acacac}
\definecolor{quarto-callout-note-color-frame}{HTML}{4582ec}
\definecolor{quarto-callout-important-color-frame}{HTML}{d9534f}
\definecolor{quarto-callout-warning-color-frame}{HTML}{f0ad4e}
\definecolor{quarto-callout-tip-color-frame}{HTML}{02b875}
\definecolor{quarto-callout-caution-color-frame}{HTML}{fd7e14}
\makeatother
\makeatletter
\@ifpackageloaded{bookmark}{}{\usepackage{bookmark}}
\makeatother
\makeatletter
\@ifpackageloaded{caption}{}{\usepackage{caption}}
\AtBeginDocument{%
\ifdefined\contentsname
  \renewcommand*\contentsname{Table of contents}
\else
  \newcommand\contentsname{Table of contents}
\fi
\ifdefined\listfigurename
  \renewcommand*\listfigurename{List of Figures}
\else
  \newcommand\listfigurename{List of Figures}
\fi
\ifdefined\listtablename
  \renewcommand*\listtablename{List of Tables}
\else
  \newcommand\listtablename{List of Tables}
\fi
\ifdefined\figurename
  \renewcommand*\figurename{Figure}
\else
  \newcommand\figurename{Figure}
\fi
\ifdefined\tablename
  \renewcommand*\tablename{Table}
\else
  \newcommand\tablename{Table}
\fi
}
\@ifpackageloaded{float}{}{\usepackage{float}}
\floatstyle{ruled}
\@ifundefined{c@chapter}{\newfloat{codelisting}{h}{lop}}{\newfloat{codelisting}{h}{lop}[chapter]}
\floatname{codelisting}{Listing}
\newcommand*\listoflistings{\listof{codelisting}{List of Listings}}
\makeatother
\makeatletter
\makeatother
\makeatletter
\@ifpackageloaded{caption}{}{\usepackage{caption}}
\@ifpackageloaded{subcaption}{}{\usepackage{subcaption}}
\makeatother
\usepackage{bookmark}
\IfFileExists{xurl.sty}{\usepackage{xurl}}{} % add URL line breaks if available
\urlstyle{same}
% fallback for those not using the hyperref driver hyperxmp:
\makeatletter
\@ifundefined{xmpquote}{\newcommand{\xmpquote}[1]{#1}}{}
\makeatother
\hypersetup{
  pdftitle={Introduction to the Free Energy Principle},
  pdfauthor={Omer Tzuk},
  colorlinks=true,
  linkcolor={blue},
  filecolor={Maroon},
  citecolor={Blue},
  urlcolor={Blue},
  pdfcreator={LaTeX via pandoc}}


\title{Introduction to the Free Energy Principle}
\author{Omer Tzuk}
\date{2026-09-05}
\begin{document}
\maketitle

\renewcommand*\contentsname{Table of contents}
{
\setcounter{tocdepth}{2}
\tableofcontents
}

\bookmarksetup{startatroot}

\chapter*{Preface}\label{preface}
\addcontentsline{toc}{chapter}{Preface}

\markboth{Preface}{Preface}

This is a Quarto book.

To learn more about Quarto books visit
\url{https://quarto.org/docs/books}.

\bookmarksetup{startatroot}

\chapter{Introduction}\label{introduction}

This is a book created from markdown and executable code.

See Knuth (1984) for additional discussion of literate programming.

\bookmarksetup{startatroot}

\chapter{What does `variational'
mean?}\label{what-does-variational-mean}

From varying a trajectory to varying a probability distribution

\hfill\break

\begin{tcolorbox}[enhanced jigsaw, arc=.35mm, bottomrule=.15mm, bottomtitle=1mm, breakable, colback=white, colbacktitle=quarto-callout-note-color!10!white, colframe=quarto-callout-note-color-frame, coltitle=black, left=2mm, leftrule=.75mm, opacityback=0, opacitybacktitle=0.6, rightrule=.15mm, title=\textcolor{quarto-callout-note-color}{\faInfo}\hspace{0.5em}{The question guiding this chapter}, titlerule=0mm, toprule=.15mm, toptitle=1mm]

In mechanics, we can vary a \emph{trajectory} and ask which trajectory
makes an action stationary. In variational inference, we vary a
\emph{probability distribution} and ask which distribution minimizes an
objective.

\textbf{``Variational'' identifies what we are doing: optimizing over
functions, here probability densities. It does not mean ``variance.''}

\end{tcolorbox}

This chapter assumes familiarity with derivatives, integrals, Gaussian
distributions, and the basic idea of conditional probability. No
previous knowledge of functional derivatives is required. The optional
derivation near the end introduces them when they become useful.

We will move through sixteen small steps. Our running example is a noisy
thermometer connected to a controller: the sensor reports a temperature,
while the actual room temperature remains uncertain. We first study
inference with the observation held fixed; controlling the heater comes
later.

By the end, you should be able to explain what is varied, derive the
free-energy identity, and distinguish an external state from a
distribution over that state.

\section{1. Start with ordinary
optimization}\label{start-with-ordinary-optimization}

Suppose we want to minimize

\[
f(x)=(x-3)^2.
\]

We vary the number \(x\). Differentiating gives

\[
f'(x)=2(x-3)=0
\quad\Longrightarrow\quad x^*=3.
\]

Because \(f''(x)=2>0\), this stationary point is a minimum. We write

\[
x^*=\arg\min_x f(x).
\]

The notation \(\arg\min\) asks for the \emph{input} that minimizes a
quantity. Here the minimizing input is \(3\), whereas the minimum value
of the function is \(f(3)=0\).

Now change the question. Instead of choosing one number, suppose we must
choose a whole curve joining two fixed points. Which curve is shortest?

For a curve written as \(y(x)\) between \(x=a\) and \(x=b\), its length
is

\[
L[y]=\int_a^b \sqrt{1+[y'(x)]^2}\,dx.
\]

Each candidate curve has one length. Thus \(L\) takes a
\textbf{function} as input and returns a \textbf{number}. Such an object
is called a \emph{functional}.

{\def\LTcaptype{none} % do not increment counter
\begin{longtable}[]{@{}
  >{\raggedright\arraybackslash}p{(\linewidth - 4\tabcolsep) * \real{0.3333}}
  >{\raggedright\arraybackslash}p{(\linewidth - 4\tabcolsep) * \real{0.3333}}
  >{\raggedright\arraybackslash}p{(\linewidth - 4\tabcolsep) * \real{0.3333}}@{}}
\toprule\noalign{}
\begin{minipage}[b]{\linewidth}\raggedright
Problem
\end{minipage} & \begin{minipage}[b]{\linewidth}\raggedright
Adjustable object
\end{minipage} & \begin{minipage}[b]{\linewidth}\raggedright
Objective
\end{minipage} \\
\midrule\noalign{}
\endhead
\bottomrule\noalign{}
\endlastfoot
Ordinary optimization & A number \(x\) & A function \(f(x)\) \\
Shortest-curve problem & A whole curve \(y(x)\) & A functional
\(L[y]\) \\
Variational inference & A probability density \(q(\eta)\) & A functional
\(F[q]\) \\
\end{longtable}
}

Square brackets help signal that the input is an entire function. They
are a notational convention, not a new mathematical operation.

\section{2. What is a variation?}\label{what-is-a-variation}

Take a candidate curve \(y(x)\) and deform it slightly:

\[
y_\epsilon(x)=y(x)+\epsilon h(x).
\]

Here \(h(x)\) describes the \emph{shape} of the deformation, and the
small number \(\epsilon\) controls its size. To preserve the endpoints,
require \(h(a)=h(b)=0\).

For each deformation shape, we ask how the length changes:

\[
\left.\frac{d}{d\epsilon}L[y+\epsilon h]\right|_{\epsilon=0}.
\]

This is the \textbf{first variation} in the direction \(h\). At a smooth
unconstrained minimum, it vanishes for every admissible direction: there
is no first-order improvement available.

The same idea appears in classical mechanics:

\[
S[x]=\int_{t_0}^{t_1}\mathcal L(x,\dot x,t)\,dt,
\qquad \delta S=0.
\]

The Euler--Lagrange equations follow by requiring stationarity under
small changes of the path with fixed endpoints. This calculus of
variations is the mathematical origin of the terminology we are
learning.

\begin{tcolorbox}[enhanced jigsaw, arc=.35mm, bottomrule=.15mm, bottomtitle=1mm, breakable, colback=white, colbacktitle=quarto-callout-important-color!10!white, colframe=quarto-callout-important-color-frame, coltitle=black, left=2mm, leftrule=.75mm, opacityback=0, opacitybacktitle=0.6, rightrule=.15mm, title=\textcolor{quarto-callout-important-color}{\faExclamation}\hspace{0.5em}{Stationary does not always mean minimum}, titlerule=0mm, toprule=.15mm, toptitle=1mm]

The condition \(\delta S=0\) means the first variation vanishes. A
stationary action need not be a minimum; it can be a saddle. We will
establish a minimum for variational free energy using the non-negativity
of KL divergence.

\end{tcolorbox}

\section{3. Enter Bayesian inference}\label{enter-bayesian-inference}

Our thermometer displays

\[
s=17^\circ\mathrm C.
\]

The true temperature \(\eta\) is hidden. Because the measurement is
noisy, the observation does not imply \(\eta=17^\circ\mathrm C\)
exactly.

A \textbf{generative model} specifies how hidden states and observations
are related:

\[
p(s,\eta)=p(s\mid\eta)p(\eta).
\]

The prior \(p(\eta)\) describes temperature uncertainty before
incorporating this reading. The likelihood \(p(s\mid\eta)\) describes
the sensor: if the room had temperature \(\eta\), how probable would
different readings be?

Bayes' rule gives the posterior:

\[
p(\eta\mid s)=\frac{p(s,\eta)}{p(s)},
\qquad
p(s)=\int p(s,\eta)\,d\eta.
\]

The denominator is called the \textbf{model evidence} or
\textbf{marginal likelihood}. It normalizes the posterior so that its
integral over \(\eta\) equals one.

For one Gaussian temperature variable, this integral is easy. For many
interacting hidden variables in a nonlinear model, it can become
prohibitively difficult. Our simple thermometer will let us learn the
method in a setting where we can check the answer exactly.

\section{4. Introduce an adjustable candidate
distribution}\label{introduce-an-adjustable-candidate-distribution}

Rather than directly calculating the posterior, introduce a candidate
density

\[
q(\eta).
\]

We adjust it to approximate \(p(\eta\mid s)\). This candidate is the
\textbf{variational density}.

Think of \(q\) as a flexible probability curve. We can move its peak,
change its width, or---if our chosen family permits---change its shape
more substantially.

Unlike an arbitrary curve, a probability density must obey

\[
q(\eta)\geq 0,
\qquad
\int q(\eta)\,d\eta=1.
\]

A small variation therefore has the form

\[
q_\epsilon(\eta)=q(\eta)+\epsilon h(\eta),
\qquad \int h(\eta)\,d\eta=0,
\]

with \(\epsilon\) and \(h\) also chosen to preserve non-negativity. We
can redistribute probability, but we cannot create or destroy its total
mass.

\section{5. What exactly is being
varied?}\label{what-exactly-is-being-varied}

The distinction is subtle but essential:

\[
\underbrace{\eta}_{\text{possible temperature}}
\qquad\text{versus}\qquad
\underbrace{q(\eta)}_{\text{density over possible temperatures}}.
\]

When we integrate over \(\eta\), we scan the possible temperatures
within \emph{one} distribution. When we vary \(q\), we change how
probability is allocated across those temperatures.

For example, changing from a broad distribution near
\(22^\circ\mathrm C\) to a narrow distribution near
\(18^\circ\mathrm C\) changes both an estimate and its uncertainty. It
does not itself cool the room.

\begin{tcolorbox}[enhanced jigsaw, arc=.35mm, bottomrule=.15mm, bottomtitle=1mm, breakable, colback=white, colbacktitle=quarto-callout-tip-color!10!white, colframe=quarto-callout-tip-color-frame, coltitle=black, left=2mm, leftrule=.75mm, opacityback=0, opacitybacktitle=0.6, rightrule=.15mm, title=\textcolor{quarto-callout-tip-color}{\faLightbulb}\hspace{0.5em}{Pause and check}, titlerule=0mm, toprule=.15mm, toptitle=1mm]

If the room stays at the same temperature while the controller becomes
more certain about it, what changes?

The physical temperature \(\eta\) can stay fixed while the density
\(q(\eta)\) becomes narrower. Inference changes a representation of the
world.

\end{tcolorbox}

\section{6. In practice, choose a family of
distributions}\label{in-practice-choose-a-family-of-distributions}

Optimizing over every conceivable density is usually impractical.
Instead, choose a family \(\mathcal Q\).

For example,

\[
q_\theta(\eta)=\mathcal N(\eta;m,v),
\qquad \theta=(m,v),\quad v>0.
\]

We use \(m\) for the mean and \(v\) for the \textbf{variance}, so the
standard deviation is \(\sqrt v\). Each pair \((m,v)\) specifies an
entire Gaussian density.

Changing \(m\) slides the curve. Changing \(v\) changes its width. A
computer can therefore optimize a two-dimensional parameter vector
rather than an unrestricted function.

\[
\min_{q\in\mathcal Q}F[q]
\quad\longleftrightarrow\quad
\min_{m,\,v>0}F[q_{m,v}].
\]

This remains variational inference: the parameters select the candidate
density. Ordinary parameter optimization is how we implement the
variational problem.

The choice of family matters. A single Gaussian cannot exactly represent
a posterior with two separated peaks. Even perfect optimization cannot
remove that limitation. This family-based approach is central to the
account in \href{https://arxiv.org/abs/1601.00670}{Blei, Kucukelbir, and
McAuliffe's review}.

\section{7. Connect the parameters to internal
states}\label{connect-the-parameters-to-internal-states}

In Free Energy Principle literature, you will often encounter

\[
q_\mu(\eta),
\]

where \(\mu\) denotes internal states that parameterize a density over
external states \(\eta\).

Here \(\mu\) need not be a scalar mean. It can be a vector whose
components specify several distribution parameters, possibly through a
nonlinear mapping \(\theta=\theta(\mu)\).

For example, a model could use two internal coordinates:

\[
\mu=(\mu_1,\mu_2),
\qquad m=\mu_1,
\qquad v=e^{\mu_2}.
\]

The exponential ensures positive variance. A physical change in \(\mu\)
then corresponds, under this specified mapping, to a change in the
candidate distribution.

\begin{tcolorbox}[enhanced jigsaw, arc=.35mm, bottomrule=.15mm, bottomtitle=1mm, breakable, colback=white, colbacktitle=quarto-callout-note-color!10!white, colframe=quarto-callout-note-color-frame, coltitle=black, left=2mm, leftrule=.75mm, opacityback=0, opacitybacktitle=0.6, rightrule=.15mm, title=\textcolor{quarto-callout-note-color}{\faInfo}\hspace{0.5em}{A mathematical representation needs justification}, titlerule=0mm, toprule=.15mm, toptitle=1mm]

Writing \(q_\mu\) defines a way to represent distributions using
internal states. It does not by itself prove that an arbitrary
thermostat performs Bayesian inference. To make that claim, we must also
justify the representation and its relation to the system's dynamics.
Neither the notation nor the existence of a Markov blanket alone
establishes posterior matching.

\end{tcolorbox}

\section{8. Give the candidate a score: KL
divergence}\label{give-the-candidate-a-score-kl-divergence}

We need a measure of how well \(q\) approximates the posterior. A
standard choice is

\[
D_{\mathrm{KL}}\!\left(q\,\Vert\,p(\cdot\mid s)\right)
=\int q(\eta)\ln\frac{q(\eta)}{p(\eta\mid s)}\,d\eta.
\]

This is an average log density ratio, with the averaging performed under
\(q\). It satisfies

\[
D_{\mathrm{KL}}\geq 0,
\]

with equality precisely when the two distributions agree almost
everywhere. ``Almost everywhere'' permits differences on sets of zero
measure, which do not change probabilities.

KL divergence is not a geometric distance: it is generally asymmetric.
The direction used here, \(D_{\mathrm{KL}}(q\Vert p)\), matters.

We would therefore like to solve

\[
q^*=\arg\min_{q\in\mathcal Q}
D_{\mathrm{KL}}\!\left(q\Vert p(\cdot\mid s)\right).
\]

Assume throughout that \(p(s)>0\) and that the relevant integrals are
finite. If \(q\) assigns positive probability where the posterior is
zero, this KL divergence is infinite.

\section{9. Remove the difficult posterior
normalization}\label{remove-the-difficult-posterior-normalization}

There is an apparent obstacle: our score contains the posterior we were
trying to avoid computing.

Use Bayes' rule inside the logarithm:

\[
\ln p(\eta\mid s)=\ln p(s,\eta)-\ln p(s).
\]

Substituting gives

\[
\begin{aligned}
D_{\mathrm{KL}}\!\left(q\Vert p(\cdot\mid s)\right)
&=\int q(\eta)
\left[\ln q(\eta)-\ln p(s,\eta)+\ln p(s)\right]d\eta\\
&=\int q(\eta)\left[\ln q(\eta)-\ln p(s,\eta)\right]d\eta
+\ln p(s)\underbrace{\int q(\eta)d\eta}_{1}.
\end{aligned}
\]

Define \textbf{variational free energy} by

\begin{equation}\protect\phantomsection\label{eq-vfe-definition}{
\boxed{
F[q;s]=\int q(\eta)\left[\ln q(\eta)-\ln p(s,\eta)\right]d\eta.
}
}\end{equation}

The semicolon emphasizes that \(s\) is held fixed while \(q\) varies. We
will sometimes abbreviate this as \(F[q]\).

Using expectation notation, the same definition is

\[
F[q;s]=\mathbb E_q\!\left[\ln q(\eta)-\ln p(s,\eta)\right].
\]

We have removed the posterior's evidence denominator from the objective.
Evaluating or estimating expectations under \(q\) may still require
work; the transformation does not make every inference problem easy.

\section{10. Understand the key
identity}\label{understand-the-key-identity}

Rearranging the previous calculation yields

\begin{equation}\protect\phantomsection\label{eq-vfe-identity}{
\boxed{
F[q;s]
=\underbrace{D_{\mathrm{KL}}\!\left(q\Vert p(\cdot\mid s)\right)}_{\text{approximation gap}}
+\underbrace{[-\ln p(s)]}_{\text{surprisal of the fixed observation}}.
}
}\end{equation}

For fixed \(s\) and a fixed model, \(-\ln p(s)\) is a constant with
respect to \(q\). Therefore

\[
\arg\min_{q\in\mathcal Q}F[q;s]
=\arg\min_{q\in\mathcal Q}D_{\mathrm{KL}}\!\left(q\Vert p(\cdot\mid s)\right).
\]

This is the payoff: \textbf{we can optimize an objective built from the
joint model and our candidate density, and thereby improve the posterior
approximation.}

Non-negativity of KL also gives

\[
F[q;s]\geq-\ln p(s).
\]

Free energy is an \textbf{upper bound} on surprisal. The gap closes when
\(q\) equals the posterior. If \(\mathcal Q\) excludes the posterior,
even its best member leaves a positive gap. A numerical optimizer can
also stop short of the best member.

Machine-learning texts often maximize the \textbf{evidence lower bound},
or ELBO:

\[
\operatorname{ELBO}[q;s]=-F[q;s]\leq\ln p(s).
\]

These are the same optimization with opposite signs.

\begin{tcolorbox}[enhanced jigsaw, arc=.35mm, bottomrule=.15mm, bottomtitle=1mm, breakable, colback=white, colbacktitle=quarto-callout-important-color!10!white, colframe=quarto-callout-important-color-frame, coltitle=black, left=2mm, leftrule=.75mm, opacityback=0, opacitybacktitle=0.6, rightrule=.15mm, title=\textcolor{quarto-callout-important-color}{\faExclamation}\hspace{0.5em}{What does inference improve?}, titlerule=0mm, toprule=.15mm, toptitle=1mm]

During this optimization, the observation and model are fixed. Improving
\(q\) decreases the approximation gap; it does \textbf{not} increase
\(p(s)\) or make the recorded observation less surprising. Model
learning and action can change other quantities, but they are different
optimization problems.

\end{tcolorbox}

\begin{tcolorbox}[enhanced jigsaw, arc=.35mm, bottomrule=.15mm, bottomtitle=1mm, breakable, colback=white, colbacktitle=quarto-callout-note-color!10!white, colframe=quarto-callout-note-color-frame, coltitle=black, left=2mm, leftrule=.75mm, opacityback=0, opacitybacktitle=0.6, rightrule=.15mm, title=\textcolor{quarto-callout-note-color}{\faInfo}\hspace{0.5em}{A note on continuous densities and units}, titlerule=0mm, toprule=.15mm, toptitle=1mm]

We use natural logarithms, giving information quantities in nats. For
continuous temperatures, \(p(s)\) is a density, not the probability of
one exact real number. Density-based surprisal depends on the
measurement coordinates and can be negative. Throughout, we use fixed
coordinates and reference measures---temperature expressed numerically
in degrees Celsius in the example. KL divergence between distributions
is invariant under a common invertible coordinate transformation.

\end{tcolorbox}

\section{11. Why call it ``free
energy''?}\label{why-call-it-free-energy}

Introduce the entropy of \(q\):

\[
H[q]=-\mathbb E_q[\ln q(\eta)].
\]

For a continuous density this is \emph{differential entropy}. Our
definition becomes

\[
F[q;s]=\underbrace{\mathbb E_q[-\ln p(s,\eta)]}_{\text{expected energy-like quantity}}-H[q].
\]

This resembles the Helmholtz free energy of thermodynamics,

\[
A=U-TS.
\]

We can make the correspondence explicit for a canonical ensemble.
Suppose physical microstates \(\eta\) have energy \(E(\eta)\) and
equilibrium density

\[
p_{\mathrm{eq}}(\eta)=\frac{e^{-\beta E(\eta)}}{Z},
\qquad\beta=\frac{1}{k_B T}.
\]

For an admissible trial density \(q\), define

\[
\mathcal A[q]=\mathbb E_q[E]-T k_B H[q].
\]

Substituting \(\ln p_{\mathrm{eq}}=-\beta E-\ln Z\) into KL divergence
gives

\[
\begin{aligned}
k_B T D_{\mathrm{KL}}(q\Vert p_{\mathrm{eq}})
&=\mathbb E_q[E]-k_B T H[q]+k_B T\ln Z\\
&=\mathcal A[q]-A_{\mathrm{eq}},
\end{aligned}
\]

where \(A_{\mathrm{eq}}=-k_B T\ln Z\). Thus the equilibrium density
minimizes this thermodynamic free-energy functional, just as the
posterior minimizes variational free energy.

The shared mathematical structure explains the terminology. However,
\(-\ln p(s,\eta)\) in an inference model is generally not a physical
Hamiltonian. \textbf{Variational free energy is not automatically an
amount of energy in joules stored in a brain or controller.}

\section{\texorpdfstring{12. Why specifically \emph{variational} free
energy?}{12. Why specifically variational free energy?}}\label{why-specifically-variational-free-energy}

We can now unpack the name precisely:

{\def\LTcaptype{none} % do not increment counter
\begin{longtable}[]{@{}
  >{\raggedright\arraybackslash}p{(\linewidth - 2\tabcolsep) * \real{0.5000}}
  >{\raggedright\arraybackslash}p{(\linewidth - 2\tabcolsep) * \real{0.5000}}@{}}
\toprule\noalign{}
\begin{minipage}[b]{\linewidth}\raggedright
Term
\end{minipage} & \begin{minipage}[b]{\linewidth}\raggedright
Meaning
\end{minipage} \\
\midrule\noalign{}
\endhead
\bottomrule\noalign{}
\endlastfoot
Variational density \(q\) & The adjustable probability distribution \\
Variational free energy \(F[q;s]\) & The functional used to score it \\
Variational inference & Optimizing that functional over an allowed
family \\
\end{longtable}
}

A second decomposition is useful in Active Inference. Factor
\(p(s,\eta)=p(s\mid\eta)p(\eta)\) in Equation~\ref{eq-vfe-definition}:

\[
\begin{aligned}
F[q;s]
&=\mathbb E_q[\ln q(\eta)-\ln p(\eta)-\ln p(s\mid\eta)]\\
&=\underbrace{D_{\mathrm{KL}}(q\Vert p(\eta))}_{\text{complexity}}
-\underbrace{\mathbb E_q[\ln p(s\mid\eta)]}_{\text{accuracy}}.
\end{aligned}
\]

Accuracy rewards candidate states that explain the observation.
Complexity measures departure from the \textbf{prior}, not necessarily
from the previous time step's belief. These terms jointly produce
posterior inference; the complexity term is not an arbitrary prohibition
on changing one's mind.

\section{13. Work through the thermometer
numerically}\label{work-through-the-thermometer-numerically}

Let temperatures be represented by their numerical values in degrees
Celsius. Choose

\[
p(\eta)=\mathcal N(\eta;m_0,v_0),
\qquad m_0=22,\quad v_0=4,
\]

and a sensor model

\[
p(s\mid\eta)=\mathcal N(s;\eta,r),
\qquad r=1,\quad s=17.
\]

The prior standard deviation is \(2^\circ\mathrm C\) and the
sensor-noise standard deviation is \(1^\circ\mathrm C\). The variances
have units of squared degrees Celsius.

Choose the variational family

\[
q(\eta)=\mathcal N(\eta;m,v).
\]

\subsection{First compute the
objective}\label{first-compute-the-objective}

For a Gaussian, \(H[q]=\tfrac12\ln(2\pi e v)\). Also,

\[
\mathbb E_q[(s-\eta)^2]=(s-m)^2+v,
\]

because the expected squared deviation includes both the squared
displacement of the mean and the variance. Similarly,

\[
\mathbb E_q[(\eta-m_0)^2]=(m-m_0)^2+v.
\]

Substitution into the free-energy definition gives

\begin{equation}\protect\phantomsection\label{eq-vfe-gaussian}{
\boxed{
\begin{aligned}
F(m,v;s)={}&\frac12\ln(2\pi r)
+\frac{(s-m)^2+v}{2r}\\
&+\frac12\ln(2\pi v_0)
+\frac{(m-m_0)^2+v}{2v_0}\\
&-\frac12\ln(2\pi e v).
\end{aligned}
}
}\end{equation}

The first line comes from the likelihood, the second from the prior, and
the last from negative entropy. Keeping the constants lets us check the
absolute free-energy bound, not only its gradients.

\subsection{Then vary the mean and
variance}\label{then-vary-the-mean-and-variance}

Differentiate with respect to \(m\):

\[
\frac{\partial F}{\partial m}
=\frac{m-s}{r}+\frac{m-m_0}{v_0}=0.
\]

Solving gives

\[
m^*=\frac{s/r+m_0/v_0}{1/r+1/v_0}=18.
\]

This is a \textbf{precision-weighted average}: precision is inverse
variance. The more precise sensor receives greater weight than the
prior.

For the variance,

\[
\frac{\partial F}{\partial v}
=\frac{1}{2r}+\frac{1}{2v_0}-\frac{1}{2v}=0,
\]

so

\[
v^*=\left(\frac1r+\frac1{v_0}\right)^{-1}=0.8.
\]

The second derivatives are positive, and there is no mixed derivative:

\[
\frac{\partial^2F}{\partial m^2}=\frac1r+\frac1{v_0}>0,
\qquad
\frac{\partial^2F}{\partial v^2}=\frac{1}{2v^2}>0.
\]

These stationary values give the unique minimum for \(v>0\).

The posterior of this Gaussian model is exactly

\[
p(\eta\mid s)=\mathcal N(\eta;18,0.8).
\]

Because our family contains it, the optimum reaches the posterior.
Notice that the answer is centered at \textbf{18}, not at the sensor
reading \textbf{17}: inference combines prior information with the
observation.

\subsection{Watch the approximation
improve}\label{watch-the-approximation-improve}

The evidence is \(p(s)=\mathcal N(s;m_0,v_0+r)\), giving

\[
-\ln p(s)=\frac12\ln(2\pi\cdot5)+\frac{(17-22)^2}{2\cdot5}
\approx4.224.
\]

The following candidates illustrate simultaneous changes in mean and
uncertainty. They are illustrative choices, not the output of a
specified learning algorithm.

{\def\LTcaptype{none} % do not increment counter
\begin{longtable}[]{@{}lrrrr@{}}
\toprule\noalign{}
Candidate & Mean \(m\) & Variance \(v\) & \(F[q;s]\) & Posterior KL
gap \\
\midrule\noalign{}
\endhead
\bottomrule\noalign{}
\endlastfoot
Initial, equal to prior & 22 & 4 & 15.419 & 11.195 \\
Intermediate & 20 & 2 & 7.016 & 2.792 \\
Near posterior & 18.5 & 1 & 4.393 & 0.170 \\
Exact posterior & 18 & 0.8 & 4.224 & 0 \\
\end{longtable}
}

Every row satisfies \(F=-\ln p(s)+D_{\mathrm{KL}}\). The evidence is
fixed throughout.

\begin{figure}

\centering{

\begin{Shaded}
\begin{Highlighting}[]
\ImportTok{import}\NormalTok{ numpy }\ImportTok{as}\NormalTok{ np}
\ImportTok{import}\NormalTok{ matplotlib.pyplot }\ImportTok{as}\NormalTok{ plt}

\NormalTok{prior\_mean, prior\_var }\OperatorTok{=} \FloatTok{22.0}\NormalTok{, }\FloatTok{4.0}
\NormalTok{observation, noise\_var }\OperatorTok{=} \FloatTok{17.0}\NormalTok{, }\FloatTok{1.0}
\NormalTok{posterior\_var }\OperatorTok{=} \FloatTok{1.0} \OperatorTok{/}\NormalTok{ (}\FloatTok{1.0} \OperatorTok{/}\NormalTok{ noise\_var }\OperatorTok{+} \FloatTok{1.0} \OperatorTok{/}\NormalTok{ prior\_var)}
\NormalTok{posterior\_mean }\OperatorTok{=}\NormalTok{ posterior\_var }\OperatorTok{*}\NormalTok{ (}
\NormalTok{    observation }\OperatorTok{/}\NormalTok{ noise\_var }\OperatorTok{+}\NormalTok{ prior\_mean }\OperatorTok{/}\NormalTok{ prior\_var}
\NormalTok{)}

\KeywordTok{def}\NormalTok{ normal\_pdf(x, mean, variance):}
    \CommentTok{"""Evaluate a normalized Gaussian density with the given variance."""}
    \ControlFlowTok{return}\NormalTok{ np.exp(}\OperatorTok{{-}}\FloatTok{0.5} \OperatorTok{*}\NormalTok{ (x }\OperatorTok{{-}}\NormalTok{ mean)}\OperatorTok{**}\DecValTok{2} \OperatorTok{/}\NormalTok{ variance) }\OperatorTok{/}\NormalTok{ np.sqrt(}
        \FloatTok{2.0} \OperatorTok{*}\NormalTok{ np.pi }\OperatorTok{*}\NormalTok{ variance}
\NormalTok{    )}

\KeywordTok{def}\NormalTok{ free\_energy(mean, variance):}
    \CommentTok{"""Evaluate the full Gaussian variational objective for fixed data."""}
\NormalTok{    likelihood\_term }\OperatorTok{=} \FloatTok{0.5} \OperatorTok{*}\NormalTok{ np.log(}\FloatTok{2.0} \OperatorTok{*}\NormalTok{ np.pi }\OperatorTok{*}\NormalTok{ noise\_var)}
\NormalTok{    likelihood\_term }\OperatorTok{+=}\NormalTok{ ((observation }\OperatorTok{{-}}\NormalTok{ mean)}\OperatorTok{**}\DecValTok{2} \OperatorTok{+}\NormalTok{ variance) }\OperatorTok{/}\NormalTok{ (}\DecValTok{2} \OperatorTok{*}\NormalTok{ noise\_var)}
\NormalTok{    prior\_term }\OperatorTok{=} \FloatTok{0.5} \OperatorTok{*}\NormalTok{ np.log(}\FloatTok{2.0} \OperatorTok{*}\NormalTok{ np.pi }\OperatorTok{*}\NormalTok{ prior\_var)}
\NormalTok{    prior\_term }\OperatorTok{+=}\NormalTok{ ((mean }\OperatorTok{{-}}\NormalTok{ prior\_mean)}\OperatorTok{**}\DecValTok{2} \OperatorTok{+}\NormalTok{ variance) }\OperatorTok{/}\NormalTok{ (}\DecValTok{2} \OperatorTok{*}\NormalTok{ prior\_var)}
\NormalTok{    entropy }\OperatorTok{=} \FloatTok{0.5} \OperatorTok{*}\NormalTok{ np.log(}\FloatTok{2.0} \OperatorTok{*}\NormalTok{ np.pi }\OperatorTok{*}\NormalTok{ np.e }\OperatorTok{*}\NormalTok{ variance)}
    \ControlFlowTok{return}\NormalTok{ likelihood\_term }\OperatorTok{+}\NormalTok{ prior\_term }\OperatorTok{{-}}\NormalTok{ entropy}

\NormalTok{surprisal }\OperatorTok{=} \FloatTok{0.5} \OperatorTok{*}\NormalTok{ np.log(}\DecValTok{2} \OperatorTok{*}\NormalTok{ np.pi }\OperatorTok{*}\NormalTok{ (prior\_var }\OperatorTok{+}\NormalTok{ noise\_var))}
\NormalTok{surprisal }\OperatorTok{+=}\NormalTok{ (observation }\OperatorTok{{-}}\NormalTok{ prior\_mean)}\OperatorTok{**}\DecValTok{2} \OperatorTok{/}\NormalTok{ (}\DecValTok{2} \OperatorTok{*}\NormalTok{ (prior\_var }\OperatorTok{+}\NormalTok{ noise\_var))}

\NormalTok{temperatures }\OperatorTok{=}\NormalTok{ np.linspace(}\DecValTok{12}\NormalTok{, }\DecValTok{29}\NormalTok{, }\DecValTok{700}\NormalTok{)}
\NormalTok{candidates }\OperatorTok{=}\NormalTok{ [(}\DecValTok{22}\NormalTok{, }\DecValTok{4}\NormalTok{), (}\DecValTok{20}\NormalTok{, }\DecValTok{2}\NormalTok{), (}\FloatTok{18.5}\NormalTok{, }\DecValTok{1}\NormalTok{), (posterior\_mean, posterior\_var)]}
\NormalTok{colors }\OperatorTok{=}\NormalTok{ [}\StringTok{"\#9b4f00"}\NormalTok{, }\StringTok{"\#6d55a3"}\NormalTok{, }\StringTok{"\#2879a8"}\NormalTok{, }\StringTok{"\#087f5b"}\NormalTok{]}
\NormalTok{styles }\OperatorTok{=}\NormalTok{ [}\StringTok{"{-}{-}"}\NormalTok{, }\StringTok{"{-}."}\NormalTok{, }\StringTok{":"}\NormalTok{, }\StringTok{"{-}"}\NormalTok{]}

\NormalTok{fig, axes }\OperatorTok{=}\NormalTok{ plt.subplots(}\DecValTok{1}\NormalTok{, }\DecValTok{2}\NormalTok{, figsize}\OperatorTok{=}\NormalTok{(}\DecValTok{11}\NormalTok{, }\FloatTok{4.3}\NormalTok{), constrained\_layout}\OperatorTok{=}\VariableTok{True}\NormalTok{)}
\ControlFlowTok{for}\NormalTok{ (mean, variance), color, style }\KeywordTok{in} \BuiltInTok{zip}\NormalTok{(candidates, colors, styles):}
\NormalTok{    axes[}\DecValTok{0}\NormalTok{].plot(}
\NormalTok{        temperatures, normal\_pdf(temperatures, mean, variance),}
\NormalTok{        color}\OperatorTok{=}\NormalTok{color, linestyle}\OperatorTok{=}\NormalTok{style, linewidth}\OperatorTok{=}\FloatTok{2.2}\NormalTok{,}
\NormalTok{        label}\OperatorTok{=}\SpecialStringTok{f"m = }\SpecialCharTok{\{}\NormalTok{mean}\SpecialCharTok{:g\}}\SpecialStringTok{, v = }\SpecialCharTok{\{}\NormalTok{variance}\SpecialCharTok{:g\}}\SpecialStringTok{"}\NormalTok{,}
\NormalTok{    )}
\NormalTok{axes[}\DecValTok{0}\NormalTok{].axvline(observation, color}\OperatorTok{=}\StringTok{"0.5"}\NormalTok{, linestyle}\OperatorTok{=}\StringTok{":"}\NormalTok{, linewidth}\OperatorTok{=}\DecValTok{1}\NormalTok{)}
\NormalTok{axes[}\DecValTok{0}\NormalTok{].}\BuiltInTok{set}\NormalTok{(xlabel}\OperatorTok{=}\StringTok{"Possible temperature η (°C)"}\NormalTok{, ylabel}\OperatorTok{=}\StringTok{"Density (1/°C)"}\NormalTok{,}
\NormalTok{            title}\OperatorTok{=}\StringTok{"Changing a whole distribution"}\NormalTok{)}
\NormalTok{axes[}\DecValTok{0}\NormalTok{].legend(frameon}\OperatorTok{=}\VariableTok{False}\NormalTok{, fontsize}\OperatorTok{=}\DecValTok{9}\NormalTok{)}

\NormalTok{means }\OperatorTok{=}\NormalTok{ np.linspace(}\DecValTok{14}\NormalTok{, }\DecValTok{23}\NormalTok{, }\DecValTok{400}\NormalTok{)}
\NormalTok{axes[}\DecValTok{1}\NormalTok{].plot(means, free\_energy(means, posterior\_var), color}\OperatorTok{=}\StringTok{"\#2879a8"}\NormalTok{,}
\NormalTok{             linewidth}\OperatorTok{=}\FloatTok{2.5}\NormalTok{, label}\OperatorTok{=}\StringTok{"F(m, v = 0.8)"}\NormalTok{)}
\NormalTok{axes[}\DecValTok{1}\NormalTok{].axhline(surprisal, color}\OperatorTok{=}\StringTok{"\#087f5b"}\NormalTok{, linestyle}\OperatorTok{=}\StringTok{"{-}{-}"}\NormalTok{,}
\NormalTok{                label}\OperatorTok{=}\StringTok{"Fixed bound: −ln p(s)"}\NormalTok{)}
\NormalTok{axes[}\DecValTok{1}\NormalTok{].scatter([posterior\_mean], [surprisal], color}\OperatorTok{=}\StringTok{"\#087f5b"}\NormalTok{, zorder}\OperatorTok{=}\DecValTok{3}\NormalTok{)}
\NormalTok{axes[}\DecValTok{1}\NormalTok{].}\BuiltInTok{set}\NormalTok{(xlabel}\OperatorTok{=}\StringTok{"Candidate mean m (°C)"}\NormalTok{, ylabel}\OperatorTok{=}\StringTok{"Free energy (nats)"}\NormalTok{,}
\NormalTok{            title}\OperatorTok{=}\StringTok{"A slice through the objective"}\NormalTok{)}
\NormalTok{axes[}\DecValTok{1}\NormalTok{].legend(frameon}\OperatorTok{=}\VariableTok{False}\NormalTok{, fontsize}\OperatorTok{=}\DecValTok{9}\NormalTok{)}
\ControlFlowTok{for}\NormalTok{ ax }\KeywordTok{in}\NormalTok{ axes:}
\NormalTok{    ax.spines[[}\StringTok{"top"}\NormalTok{, }\StringTok{"right"}\NormalTok{]].set\_visible(}\VariableTok{False}\NormalTok{)}
\NormalTok{    ax.grid(alpha}\OperatorTok{=}\FloatTok{0.15}\NormalTok{)}
\NormalTok{plt.show()}
\end{Highlighting}
\end{Shaded}

}

\caption{\label{fig-varying-density}}

\end{figure}%

\subsection{What if the family is too
restricted?}\label{what-if-the-family-is-too-restricted}

Suppose we fix \(v=4\) and vary only the mean. The best mean is still
\(18\), but \(q=\mathcal N(18,4)\) remains too broad. Its residual KL
divergence is

\[
\frac12\left[\frac{4}{0.8}-1-\ln\frac{4}{0.8}\right]
\approx1.195\ \text{nats}.
\]

Matching the posterior mean is not the same as matching the posterior
distribution.

\section{14. Make the word ``variation''
tangible}\label{make-the-word-variation-tangible}

Imagine a flexible sheet whose shape represents a density. You adjust
the sheet and evaluate a score. If the score falls, the new shape is a
better candidate.

Unlike an unconstrained rubber sheet, this one must stay nonnegative and
enclose unit area. Its permitted shapes depend on your family
\(\mathcal Q\). A Gaussian family permits shifting and stretching a
single bell-shaped curve.

The analogy identifies three separate ingredients: a deformable object
\(q\), admissible deformations, and a scalar score \(F[q;s]\). The
following optional calculation turns that picture into calculus.

\begin{tcolorbox}[enhanced jigsaw, arc=.35mm, bottomrule=.15mm, bottomtitle=1mm, breakable, colback=white, colbacktitle=quarto-callout-note-color!10!white, colframe=quarto-callout-note-color-frame, coltitle=black, left=2mm, leftrule=.75mm, opacityback=0, opacitybacktitle=0.6, rightrule=.15mm, title=\textcolor{quarto-callout-note-color}{\faInfo}\hspace{0.5em}{Optional: minimize over the full density}, titlerule=0mm, toprule=.15mm, toptitle=1mm]

Instead of assuming Gaussian form, minimize \(F[q;s]\) over normalized
densities on the posterior support. Introduce a Lagrange multiplier
\(\lambda\) for normalization:

\[
\mathcal J[q]=F[q;s]+\lambda\left(\int q(\eta)d\eta-1\right).
\]

Perturb \(q\) to \(q+\epsilon h\). Since \(d(q\ln q)/dq=\ln q+1\), the
first variation is

\[
\left.\frac{d\mathcal J[q+\epsilon h]}{d\epsilon}\right|_{\epsilon=0}
=\int h(\eta)\left[\ln q(\eta)+1-\ln p(s,\eta)+\lambda\right]d\eta.
\]

For this to vanish for arbitrary admissible perturbations in the
interior of the support, the bracket must vanish:

\[
\frac{\delta\mathcal J}{\delta q(\eta)}
=\ln q(\eta)+1-\ln p(s,\eta)+\lambda=0.
\]

The expression \(\delta\mathcal J/\delta q(\eta)\) is the
\textbf{functional derivative}. It describes sensitivity to changing the
density at each location \(\eta\).

Rearranging,

\[
q^*(\eta)=e^{-1-\lambda}p(s,\eta).
\]

Normalization requires

\[
1=e^{-1-\lambda}\int p(s,\eta)d\eta
=e^{-1-\lambda}p(s).
\]

Therefore

\[
q^*(\eta)=\frac{p(s,\eta)}{p(s)}=p(\eta\mid s).
\]

We have recovered Bayes' rule by varying a density. The KL identity
establishes that this stationary density is a global minimum. This
derivation identifies the mathematical optimum; it does not supply a
general way to compute the evidence integral.

\end{tcolorbox}

\section{15. Return to the Free Energy
Principle}\label{return-to-the-free-energy-principle}

Read the notation \(q_\mu(\eta)\) as:

\begin{quote}
An adjustable candidate distribution over external states, parameterized
by internal states.
\end{quote}

For a designed inference algorithm, one possible update is gradient
descent:

\[
\dot\mu=-\kappa\nabla_\mu F[q_\mu;s],
\qquad\kappa>0.
\]

If \(s\) and the model stay fixed, this gives

\[
\frac{dF}{dt}
=\nabla_\mu F\cdot\dot\mu
=-\kappa\|\nabla_\mu F\|^2\leq0.
\]

That is a concrete example of internal-state dynamics implementing
variational optimization. In our Gaussian example, descending in the
mean uses

\[
\dot m=-\kappa\left[\frac{m-s}{r}+\frac{m-m_0}{v_0}\right].
\]

The two precision-weighted discrepancies pull the estimate toward the
observation and the prior mean.

This deliberately simple descent equation is not a general theorem about
every physical system. Stochastic dynamics, changing sensory inputs, and
solenoidal flow require additional treatment. In
\href{https://arxiv.org/abs/2201.06387}{Friston and colleagues'
account}, relating physical dynamics to Bayesian descriptions is a
further argument built on assumptions about the system and its
partition.

For orientation, that paper also uses particular states
\(\pi=(s,a,\mu)\). The corresponding free-energy identity can be written

\[
F[q_\mu;\pi]=D_{\mathrm{KL}}\!\left(q_\mu(\eta)\Vert p(\eta\mid\pi)\right)-\ln p(\pi).
\]

The algebra has the same form, but \(\mu\) now appears both in the
density parameters and in \(\pi\). Our fixed-observation calculation
does not by itself establish the paper's claims about gradients of
physical states. In particular, exact identities involving a specially
constructed density and conditional modes should not be read as saying
that every candidate \(q_\mu\) equals the posterior.

Finally, distinguish \textbf{inference} from \textbf{action}. Updating
\(q\) revises the temperature estimate. Turning on a heater changes the
environment and subsequent observations. To model the latter, we need a
model of how actions affect observations and how preferred outcomes or
viable states enter the controller's model. An inference prior centered
at \(22^\circ\mathrm C\) alone is not a complete heating policy.

Variational free energy is also distinct from \textbf{expected free
energy}, an objective used in many Active Inference models to evaluate
possible future policies. Understanding the present chapter provides the
foundation for that later topic.

\section{16. Keep three different objects
separate}\label{keep-three-different-objects-separate}

{\def\LTcaptype{none} % do not increment counter
\begin{longtable}[]{@{}
  >{\raggedright\arraybackslash}p{(\linewidth - 4\tabcolsep) * \real{0.3333}}
  >{\raggedright\arraybackslash}p{(\linewidth - 4\tabcolsep) * \real{0.3333}}
  >{\raggedright\arraybackslash}p{(\linewidth - 4\tabcolsep) * \real{0.3333}}@{}}
\toprule\noalign{}
\begin{minipage}[b]{\linewidth}\raggedright
Symbol
\end{minipage} & \begin{minipage}[b]{\linewidth}\raggedright
Meaning
\end{minipage} & \begin{minipage}[b]{\linewidth}\raggedright
Thermometer example
\end{minipage} \\
\midrule\noalign{}
\endhead
\bottomrule\noalign{}
\endlastfoot
\(\eta\) & The hidden state itself & Actual room temperature \\
\(p(\eta\mid s)\) & Exact posterior under the chosen model &
Distribution implied by prior, sensor model, and reading \\
\(q_\mu(\eta)\) & Candidate distribution & Distribution represented by
adjustable internal parameters \\
\end{longtable}
}

``Exact posterior'' means exact \textbf{under the model}. It does not
guarantee that the model correctly describes the real sensor or room.

The central optimization is

\[
\boxed{
\mu^*=\arg\min_\mu F[q_\mu;s].
}
\]

Varying \(\mu\) varies a distribution. At a global optimum, we obtain
the member of the allowed family with the smallest posterior KL
divergence. If the family contains the posterior, that divergence can
reach zero.

\begin{tcolorbox}[enhanced jigsaw, arc=.35mm, bottomrule=.15mm, bottomtitle=1mm, breakable, colback=white, colbacktitle=quarto-callout-tip-color!10!white, colframe=quarto-callout-tip-color-frame, coltitle=black, left=2mm, leftrule=.75mm, opacityback=0, opacitybacktitle=0.6, rightrule=.15mm, title=\textcolor{quarto-callout-tip-color}{\faLightbulb}\hspace{0.5em}{Explain it in one sentence}, titlerule=0mm, toprule=.15mm, toptitle=1mm]

\textbf{Variational inference turns posterior estimation into
optimization over probability distributions; variational free energy is
the objective whose excess above surprisal measures the remaining
posterior mismatch.}

\end{tcolorbox}

\section*{Check your understanding}\label{check-your-understanding}
\addcontentsline{toc}{section}{Check your understanding}

\markright{Check your understanding}

\begin{enumerate}
\def\labelenumi{\arabic{enumi}.}
\tightlist
\item
  In \(F[q;s]\), which object is varied and which observation is held
  fixed?
\item
  Why can free energy fall while \(p(s)\) stays unchanged?
\item
  In the Gaussian example, why is the posterior mean \(18\) rather than
  \(17\)?
\item
  If the mean is correct but the variance is too large, is inference
  exact?
\item
  Does \(\delta S=0\) always mean that the action is minimized?
\item
  If we allow an arbitrary normalized density, what density minimizes
  \(F\)?
\end{enumerate}

\begin{tcolorbox}[enhanced jigsaw, arc=.35mm, bottomrule=.15mm, bottomtitle=1mm, breakable, colback=white, colbacktitle=quarto-callout-note-color!10!white, colframe=quarto-callout-note-color-frame, coltitle=black, left=2mm, leftrule=.75mm, opacityback=0, opacitybacktitle=0.6, rightrule=.15mm, title=\textcolor{quarto-callout-note-color}{\faInfo}\hspace{0.5em}{Suggested answers}, titlerule=0mm, toprule=.15mm, toptitle=1mm]

\begin{enumerate}
\def\labelenumi{\arabic{enumi}.}
\tightlist
\item
  We vary \(q\), or the parameters selecting it. We hold \(s\) and the
  generative model fixed.
\item
  The posterior KL gap falls; the surprisal term is constant.
\item
  The estimate combines the prior mean \(22\) and reading \(17\),
  weighted by their precisions.
\item
  No.~Exact inference matches the full posterior distribution, including
  its uncertainty.
\item
  No.~It expresses stationarity of the action.
\item
  The posterior \(p(\eta\mid s)\), subject to the support and regularity
  conditions discussed above.
\end{enumerate}

\end{tcolorbox}

\section*{Further reading}\label{further-reading}
\addcontentsline{toc}{section}{Further reading}

\markright{Further reading}

\begin{itemize}
\tightlist
\item
  David M. Blei, Alp Kucukelbir, and Jon D. McAuliffe (2017),
  \href{https://arxiv.org/abs/1601.00670}{\emph{Variational Inference: A
  Review for Statisticians}}. A systematic introduction to variational
  families, KL minimization, and the ELBO.
\item
  Karl Friston, Lancelot Da Costa, Noor Sajid, Conor Heins, Kai
  Ueltzhöffer, Grigorios A. Pavliotis, and Thomas Parr (2023),
  \href{https://arxiv.org/abs/2201.06387}{\emph{The free energy
  principle made simpler but not too simple}}. The connection to random
  dynamical systems and the Bayesian interpretation of internal states;
  the preprint first appeared in 2022.
\end{itemize}

\bookmarksetup{startatroot}

\chapter{Summary}\label{summary}

In summary, this book has no content whatsoever.

\bookmarksetup{startatroot}

\chapter*{References}\label{references}
\addcontentsline{toc}{chapter}{References}

\markboth{References}{References}

\protect\phantomsection\label{refs}
\begin{CSLReferences}{1}{1}
\bibitem[\citeproctext]{ref-knuth84}
Knuth, Donald E. 1984. {``Literate Programming.''} \emph{Comput. J.}
(USA) 27 (2): 97--111. \url{https://doi.org/10.1093/comjnl/27.2.97}.

\end{CSLReferences}




\end{document}
